% SOURCE docs/release/2026-10-29/source/docs/writing-repair/2026-10-29/assistance.tex CHAR 0
\chapter{Which tasks merit language assistance?}
\label{rewrite:assistance-tasks}

Assistance should be selected by the task it improves. A system that answers
questions about a supplied causal mechanism has not thereby shown that
it can discover a defensible mechanism for a new bank. The evidence
synthesis in Section~\ref{scisyn:assistance} and the detailed studies
support a restricted comparison, with a questionnaire and ordinary search
as serious alternatives. They do not establish a measured saving for this
proposed product.



% SOURCE docs/release/2026-10-29/source/docs/writing-repair/2026-10-29/src/e2b7999e21232da9.tex CHAR 0
\section{Mechanisms, timing and the interpretation of a diagram}
\label{scisyn:assistance}
A retention request can mean an invitation, acceptance, a price change or a
repeated contact rule. The distinction determines the population and
intervention before any graph is evaluated. Question clarification and graph
critique are plausible assisted tasks in the inspected literature.
\citep{v12_la_2305_00050,v12_la_2606_23608}

Each proposed edge needs an economic interpretation, an information date and
supporting evidence. A graph with unresolved directions can already guide a
decision when every credible member implies the same action. Conversely, one
wrong orientation can reverse a financial sign even when most edges are
correct. The mixed LACR retrieval results support testing the retrieval choice
itself. \citep{v12_la_2402_15301}

Formal reasoning under a supplied mechanism and eliciting the mechanism of a
new bank are different tasks. CLadder and causal-learn+ motivate different
comparisons on that distinction. In a repeated policy, later observations
cannot silently become pretreatment adjustment variables. An unsupported
calculation calls for an explanation of the missing information or a revised
question. \citep{v12_la_2312_04350,v12_la_2606_23608}


% SOURCE docs/release/2026-10-29/source/docs/writing-repair/2026-10-29/assistance.tex CHAR 594


\section{Formulation and variable construction}
A useful formulation states the intervention and its alternative, the
population, outcome, horizon and information date, and the choices that remain
unresolved. The retained Kıcıman and causal-learn+ discussions motivate a
prototype for question clarification and graph critique; they cannot establish
the truth of a bank's assumptions. A structured questionnaire completed by an
economist is therefore a serious comparator. The relevant outcomes are material
omissions in the final question and the time needed to correct them. In the
deposit example, assignment, receipt, acceptance and price define different
questions. A fluent response that silently substitutes one for another fails
this comparison. \citep{v12_la_2305_00050,v12_la_2606_23608}

Variable construction introduces a further question: does an extracted field
measure the scientific variable? The retained review of text-based adjustment
explains relevant uses of text, but supplies no general guarantee that an
embedding removes confounding. A documented coding rule and independently
labelled examples provide a comparison for the proposed extraction, including
missing and ambiguous cases. The observation journey explains why predictive
accuracy and apparent completeness alone cannot establish measurement validity.
Auditable field definitions and retrieval assistance are justified starting
points; autonomous scientific measurement requires independent validation of
the relation between the field and the variable. \citep{exec13_census_text_review_arxiv}



% SOURCE docs/release/2026-10-29/source/docs/writing-repair/2026-10-29/src/0dc21980bfecb7cc.tex CHAR 0
\section{Measurement accuracy and causal validity}
\label{reading:assistance-conclusions}

\subsection{Better measurement can still leave confounding}

Extracting financial distress from customer notes illustrates the distinction.
A model may summarize the notes accurately while the notes omit circumstances
that influenced both a servicing action and repayment. The measurement example
in Section~\ref{exec13:census-measurement} makes the consequence explicit.
With true distress $U$, measured distress $W=U+\varepsilon_W$, action
$A=U+\varepsilon_A$ and outcome $Y=\tau A+U+\varepsilon_Y$, mutually
independent centred Gaussian disturbances yield
\[
 \beta_{A\mid W}=\tau+\frac{s}{1+2s},\qquad
 s=\operatorname{Var}(\varepsilon_W),
\]
when the other variances are one. The second term is residual confounding.
It persists in an arbitrarily large sample. A narrower statistical interval
around the adjusted coefficient can therefore make an incorrect economic
interpretation appear more certain. The primary review of text as confounder
measurement discusses the need to distinguish confounders, mediators and
colliders; the displayed calculation is the manuscript's constructed example.
\citep{exec13_census_text_review_arxiv}

\begin{figure}[htbp]\centering
\begin{tikzpicture}[x=2.2cm,y=6cm,>=Stealth,every node/.style={font=\small}]
 \draw[->] (0,-.25)--(0,.3) node[above]{Effect units};
 \draw[->] (0,-.25)--(4.35,-.25) node[right]{Measurement-noise variance};
 \draw[dashed,gray] (0,0)--(4.1,0);
 \draw[thick,ink,dashed] (0,-.2)--(4,-.2);
 \draw[very thick,teal,domain=0:4,samples=60] plot (\x, {-.2+\x/(1+2*\x)});
 \foreach \x in {1,2,3,4}{\draw (\x,-.26)--(\x,-.24);\node[below] at (\x,-.27){\x};}
 \node[left] at (0,-.2){$-0.2$};\node[left] at (0,0){$0$};\node[left] at (0,.2){$0.2$};
 \node[above,align=center] at (2.7,.23){Adjusted regression\ coefficient};
 \node[above] at (2.7,-.2){Specified causal effect};
\end{tikzpicture}
\caption{In the constructed Gaussian example, measurement error can reverse
the sign of an adjusted regression coefficient. The causal effect is fixed at
$\tau=-1/5$; the solid curve is $-1/5+s/(1+2s)$. This is an exact population
calculation, not observed performance of a language model.}
\label{reading:measurement-error}
\end{figure}

The appropriate comparison is consequently broader than extraction accuracy.
Simple dictionaries, supervised extraction, human coding and language models
should use the same permitted document dates and the same causal target.
The evaluation should ask whether their measurement errors change the estimated
effect or the chosen action. It should also test cases in which the notes were
written after treatment or only after a complaint. Those cases alter the
causal role of the text, not merely its predictive quality.
\citep{exec13_census_text_review_arxiv}

This example concerns the validity of the measurement used for adjustment.
The task comparisons in Chapter~\ref{rewrite:assistance-tasks} separately
address formulation, discovery, formal reasoning, execution and reporting.
They distinguish supplied-structure benchmarks from elicitation of a new
bank mechanism. \citep{v12_la_2312_04350,v12_la_2606_23608}


% SOURCE docs/release/2026-10-29/source/docs/writing-repair/2026-10-29/src/80a538d3e6bbabed.tex CHAR 0
\section{What repeated readings add to measurement}
\label{gap29:measurement}

A bank may ask an assistant to recover financial distress from customer
notes before assessing the effect of a servicing action. Reading the same
notes several times is attractive because disagreement is easy to observe.
Agreement can be less informative: several readers may share the same
omission. The distinction affects whether text can help control confounding.
The literature on text-based causal adjustment identifies measurement error,
collider information and post-treatment content as separate problems.
\citep{exec13_census_text_review_arxiv}

\subsection{A common error creates a lower limit}
Consider a project model, not an empirical claim about any language model.
Let latent distress be $U$, the action index $A=U+\varepsilon_A$, and
the outcome $Y=\tau A+U+\varepsilon_Y$, with $\tau=-1/5$. The latent
variable and the two disturbances are independent, centred, unit-variance
Gaussians. Suppose reading $j$ produces $W_j=U+e_j$. The measurement errors
are independent of $U,\varepsilon_A,\varepsilon_Y$, have variance $s$ and
equal pairwise correlation $\rho$. For this comparison the measurement
variance is positive and the common correlation is nonnegative and strictly
below one. These restrictions permit a nonsingular error covariance for any
number of readings. A fixed negative correlation would eventually make that
covariance invalid as the number increased. The limiting argument below
holds the admissible common correlation fixed. These premises make the
calculation transparent; actual extraction errors may violate all of them.

The averaged reading has error variance
\[
 s_m=\operatorname{Var}\!\left(\frac{1}{m}\sum_{j=1}^m e_j\right)
     =s\left(\rho+\frac{1-\rho}{m}\right).
\]
The covariance of $(A,\overline W)$ is
$\left(\begin{smallmatrix}2&1\\1&1+s_m\end{smallmatrix}\right)$,
and its covariance with $Y$ is $(2\tau+1,\tau+1)'$. Solving these two
normal equations gives the coefficient on $A$ after adjustment:
\[
 \beta_{A\mid\overline W}=\tau+\frac{s_m}{1+2s_m}.
\]
Under exchangeability, the contrasts between readings contain only noise
orthogonal to both $A$ and $Y$. Including every $W_j$ therefore yields the
same population coefficient as adjusting for their mean. Exact rational
calculations agree with the full $(m+1)$-variable covariance solution in all
20 prespecified cases.

For $\tau=-1/5$, the adjusted coefficient is negative only if $s_m<1/3$.
With independent unit-variance errors, at least four readings cross that
threshold. With error correlation $1/5$, at least seven are needed. If
$s\rho\geq 1/3$, no finite number can repair the sign. The measurement
variance tends to $s\rho$, so repeated readings retain a shared error floor.
At $s=1$ and $\rho=1/2$, the limiting adjusted coefficient is $1/20$,
although the true effect is $-1/5$.

\begin{figure}[htbp]
\centering
\includegraphics[width=.94\textwidth]{scientific-gap-closure/2026-10-29/measurement-floor.pdf}
\caption{Project-derived population coefficients under repeated noisy
measurements. Curves are exact calculations under the stated Gaussian model.
The true coefficient is -0.2; positive coefficients have the wrong sign.
No language model was evaluated.}
\label{gap29:measurement-floor}
\end{figure}

\subsection{An evaluation should test the proposed repair}
This calculation gives a reason to measure errors against an independent
reference rather than to count agreement between prompts. A validation sample
could compare a structured dictionary, a single extraction, repeated
extractions and retrieval-assisted extraction using the same permitted notes
and decision date. Each method should record omitted information, timing
violations and error dependence across cases. An outcome observed after the
action must not quietly enter the pretreatment measurement.

The economic test is whether the measured improvement changes the causal
conclusion at an acceptable cost. Lower text-classification error alone is
insufficient if the remaining error tracks the action or outcome. Conversely,
the shared-error limit suggests a concrete remedy: obtain an independent
measurement or better source information, rather than simply requesting more
readings of the same incomplete account. A second model is not automatically
an independent measurement. Independence requires evidence.

Graph construction raises the same issue. An assistant can propose
mechanisms and competing diagrams, but the relevant evidence must concern
the bank's assignment and information process. A formal method can evaluate
an intervention under a supplied graph; it cannot establish the graph's
missing arrows from fluent prose. The next prototype comparison should
therefore count consequential formulation and graph errors, alongside time
spent correcting them. It should keep a questionnaire plus ordinary search
as an alternative and preserve the formal method's unsupported conclusions.

The present derivation establishes a limit and a discriminating experiment.
It measures no live model benefit, does not close the unread literature
census and supplies no universal measurement-error correction. Its practical
contribution is to identify which additional information could make repeated
assistance useful, and which repetition cannot supply.


% SOURCE docs/release/2026-10-29/source/docs/writing-repair/2026-10-29/assistance.tex CHAR 2293


\section{Diagrams and numerical discovery}
A proposed diagram is useful when each relationship has an interpretation,
an information date and supporting evidence, and competing explanations remain
visible. LACR's retrieval results show why the choice of evidence matters: the
retained study reports both improvement and deterioration, including the ASIA
skeleton comparison and the SACHS comparison after adding PC output. More
context is therefore a candidate design choice whose consequences need
measurement. \citep{v12_la_2402_15301}

Ordinary search by an economist and an applicable numerical discovery method
are substantive alternatives. A comparison with matched inputs should follow
false orientations and missed common causes through to the estimated effect
or chosen action. Edge counts alone conceal that consequence. The first
assisted workflow can propose a diagram and its alternatives for review, then
invoke a named discovery or identification method under explicit assumptions.
Numerical discovery contributes evidence within those assumptions; textual
plausibility cannot resolve directions that the observations leave equivalent.
In the three-model pricing example, one unjustified orientation is enough to
reverse the financial sign.



% SOURCE docs/release/2026-10-29/source/docs/writing-repair/2026-10-29/src/02a318fabcd8ae09.tex CHAR 0
\section{What causal reasoning benchmarks establish}
\label{gap29:frontier}
A banking analyst asks several different questions of an assistant. Can it
translate a request for higher deposits into an intervention and a measurable
target? Can it distinguish information known before an offer from information
created by the offer? Can it explain which arrows in a diagram are assumptions,
and recognise when several diagrams fit the same observations? A correct
answer to a narrative puzzle is useful evidence about some of these abilities,
but it does not establish all of them.

\subsection{Explanations deserve their own test}
Wang and Shen examine three constructed scenarios: coin operations, interrupted
file downloads and a hypothetical vaccination calculation. Four variants of
each scenario and three prompting approaches give 36 prompt conditions.
Their scoring distinguishes a correct result and explanation, a flawed
explanation, and an incorrect result. This is valuable because a correct
number can accompany a mistaken economic argument. The study reports mixed
changes from examples and guided reasoning, rather than a uniformly helpful
prompting intervention. \citep{gap29wangshen}

The terms used in a benchmark also need scrutiny. The paper calls reasoning
from an observed effect to a cause \emph{actual causality}. That task description
does not, by itself, specify a formal criterion for an individual's actual
cause. Its smallpox case asks which source produces more deaths in a stipulated
population. With the given numbers, vaccination reactions produce
$990{,}000(0.01)(0.01)=99$ deaths, while disease among unvaccinated people
produces $10{,}000(0.02)(0.20)=40$. Those population counts differ from the
per-person risks, $0.0001$ and $0.004$. Confusing a total with a risk is an
estimand error; the arithmetic alone does not establish a general
Simpson-reversal mechanism. These are deductions from the paper's stated
example, not new medical evidence.

The banking analogue is direct. A large customer group can contribute more
defaults while offering a smaller default risk per account. A deposit offer
can increase gross balances without increasing net contribution. An assistant
should state the denominator and economic target before presenting the
answer. The paper's prompts provide useful stress cases for that distinction.
The small scenario set and model configurations do not support a current
ranking of assistants for banking work. We have inspected the published
methods and results; we have not reproduced its model calls or inspected every
supplementary response.

\subsection{Scale does not remove uncertainty about the answer key}
MDC-Bench broadens the comparison to 8,422 questions across seven disciplines,
including business, economics and management. Its tasks distinguish reasoning,
intervention, decision and counterfactual questions. Questions are constructed
from causal narratives extracted from preprints, with model-generated
alternatives and a manual sampling check. The published protocol identifies
model versions, prompt templates, temperature and output limits.
\citep[Sections 2--3 and Appendix A]{gap29mdc}

Its harder condition requires two ranked answers instead of one. For example,
the reported Claude thinking configuration has average accuracy 84.56\% in
the basic condition and 59.25\% in the harder condition. These are different
response requirements on the same benchmark; their difference is not a
measured deterioration in one fixed banking task. The benchmark is useful
for comparing errors by causal structure and question type, provided its
answer-generation assumptions remain visible.

The detailed prompts reveal an important distinction. Reverse reasoning may
use necessity and sufficiency, or otherwise select an abductive explanation
judged probable from background knowledge. Intervention questions construct
an altered causal narrative; the answer alternatives are ranked by their
judged plausibility. These procedures do not supply a known structural
equation system, a randomised treatment effect or an identified posterior
over causal diagrams. Appendix A's sampling review helps detect poor
questions but does not turn model-generated causal statements into independent
economic truth. The authors also discuss weak long-chain relations and
distractors that may be too easy. \citep[Limitations; Tables 28--29 and
39--40]{gap29mdc}

\subsection{An agent architecture still needs an identified effect}
Chen combines language models, diffusion models and interacting agents in a
study of climate-policy sentiment. Section 6 describes 43,943 tweets and
comparisons with topic, language and diffusion baselines. It reports F1,
AUC, causal-chain precision, counterfactual error and simulated intervention
outcomes. Table 4 removes individual components to examine their contribution.
These are reported results of the paper; the model was not reproduced
locally. \citep[Sections 5.6 and 6; Tables 1--4]{gap29chen}

The causal interpretation needs more than those comparisons. Equation (16)
factorises a joint distribution according to a graph. Equation (17) defines
a conditional treatment contrast. Neither statement alone identifies that
contrast from the observed tweets. Equation (19) compares predicted and
reference counterfactual outcomes, but its accompanying text allows the
reference to come from annotations or simulated baselines. Section 6.1 does
not supply enough detail to reconstruct an independently identified
counterfactual for each tweet. F1 and AUC assess labelled prediction; they
do not repair that missing identification argument.

The useful proposal lesson is to separate an architecture from the evidence
for a decision. A bank may test whether a language model helps extract a
policy variable or whether a simulator represents a declared mechanism.
It still needs a defensible treatment assignment, observation model and
counterfactual comparison before attributing a financial gain to the policy.
The reported architecture is a research lead. Its stated improvements are
not evidence that the proposed banking service already delivers causal or
commercial benefits.

\subsection{A useful banking comparison}
The proposal should use such benchmarks to diagnose formulation and reasoning
failures, then add tasks whose consequences can be checked independently.
A compact questionnaire and an executable calculation are necessary
comparators. For a deposit question, all methods should receive the same
decision date, eligibility rule, alternative offer and contribution definition.
For servicing, all should see the same observed history and reasons for
dropout. For a diagram, an admissible set should remain admissible when the
observations cannot orient an edge.

The resulting assessment has distinct outcomes: whether the target is correct,
whether a proposed adjustment is justified, whether an unknown effect remains
unknown, whether the numerical calculation is correct, and how much analyst
time the assistance saves. Independent adjudication is needed for open-ended
economic formulations. Generated cases with known mechanisms can test
specific errors before an expensive bank exercise. The public experimental,
servicing and linkage studies in this volume provide numerical cases of that
kind, but no live assistant has yet been evaluated on them.

The inspected papers therefore sharpen the proposed experiment rather
than establishing a benefit from the proposed service. Repeated agreement
among prompts is especially weak evidence when their errors share a source,
as the repeated-measurement calculation shows. The useful next comparison
measures consequential mistakes and time on held-out tasks, with human
adjudication and a fixed cost allowance.


% SOURCE docs/release/2026-10-29/source/docs/writing-repair/2026-10-29/assistance.tex CHAR 3610


\section{A positive workflow result and its limits}
The retained ORCA study gives a more direct reason to test assistance in a
completed workflow. Twelve participants used a counterbalanced within-subjects
design; the baseline assistant already had tools, step gates and minimal
state. The paper's Figure 5 reports total task durations of 1,395 seconds
for the baseline and 1,107 seconds for ORCA, a reported reduction of about
21\%. Table 1 reports completion rates 0.75 and 0.83. This is positive
published evidence for a structured workflow on the study's relational,
generated causal tasks. It is not a measured benefit for this banking
proposal. \citep[pp.~3--5]{exec13_census_orca}

The interpretation changes when the outcome definitions and all attempted
analyses are retained. Completion means producing an ATE, regardless of
accuracy; discovery success requires the treatment and outcome nodes to be
present. The paper's confidence-interval metric is described as an estimated
ATE falling within a ground-truth interval, so its reported 0.083 and 0.50
must not be presented as conventional repeated-sample confidence coverage.
The full-data condition comparison for absolute ATE error has $p=0.321$;
after excluding configuration-error outliers, the reported OLS comparison
has $p=0.020$. Both belong in an assessment of completed work because
configuration errors are among the failures a service is supposed to avoid.
The abstract's claimed 42-percentage-point completion gain also differs
from the displayed 0.75-to-0.83 comparison; this synthesis uses the stated
table comparison rather than reconciling that discrepancy by invention.
\citep[pp.~1, 5, 17--18]{exec13_census_orca}

The proposal should therefore test whether shared intermediate results and
question clarification save full effort while preserving scientific quality.
It should retain all configuration failures and report separately the time
spent by the analyst, tool execution and correction. The external study
supports that experiment's plausibility. Its small participant sample,
specified tasks, outcome definitions and analysis exclusions prevent a
claim that the same benefit is established across bank studies.

\section{Identification, execution and reporting}
CLadder evaluates formal causal questions under constructed mechanisms. Its
historical model results concern a supplied task rather than the elicitation
of a new bank mechanism. A defined formal procedure provides the identification
comparator: an assisted answer must preserve both the conclusion and the
conditions under which it follows. When identification fails, an explanation
of the missing restriction is useful; an unexplained empty response is not.
The benchmark literature and retained mechanism-clustering critiques also
justify retaining failed cases and the full uncertainty denominator.
\citep{v12_la_2312_04350}

For execution, the local interface study establishes equality for four
constructed requests. Assessing assistance requires a further comparison with
a scripted workflow, because choosing the request and correcting mistakes are
part of the analytical work. Request construction, faithful tool use, detected
errors and correction effort are the relevant outcomes. The formal tool's
calculation cannot by itself establish the assistant's contribution.

A deterministic reporting template is another strong comparator. It already
states the estimand, assumptions, interval and scope. The proposed additional
value is an explanation of why the result changes a decision, or why it
cannot yet support one, that remains faithful to the result and its sources.
An explanation that promotes a set crossing zero into a positive recommendation
is substantively wrong. No local live-model evaluation yet establishes an
incremental benefit for this reporting task.

Question clarification with evidence retrieval and explicit alternatives is
the preferred first assisted task, evaluated against the questionnaire.
Identification and numerical results remain with defined procedures. Expansion
to autonomous discovery, variable construction or recommendations needs
evidence specific to those tasks. The three frontier studies in
Section~\ref{gap29:frontier} sharpen this distinction. The financing-and-exports
study in Section~\ref{closure30:formulation} supports a bounded extraction and
explanation comparison; its scores do not identify a lending effect. The
related epistemology discussion supplies no separate intervention estimate.
These findings refine the proposed experiment without establishing a live
banking benefit. \citep{rel18li,rel18mugleston}

The end-to-end workflow in Section~\ref{rel64:graph-confidence} turns a
prompt-based risk score into sampled graphs and fitted effect intervals. Its
uncertainty interpretation depends on the additional conditions developed
there. The thesis by Ait Lhouss is used only to explain a source-quality
exclusion: Section~4.1 and the selected-study table report ten studies, whereas
the abstract and Tables~5--6 report eighty. These unresolved denominators
prevent using its reported review coverage or productivity results. That
limited exclusion implies neither intent nor a formal retraction.
\citep{rel64molero,rel64aitlhouss}

The pharmacovigilance comparison in Section~\ref{rel66:text-signals} shows
that useful enriched terms and classifier-score differences still need an
intervention argument. The InferBERT and DeepCausality methods expose the
missing adjustment requirements; the conditional text-embedding alternative
also makes target-population weighting explicit. The fixed frontier inventory
contains 91 inspected entries, including the one source-quality quarantine,
and 22 unread entries. Supporting method originals are audited separately
and do not enlarge that denominator. Inspection is therefore distinct from
positive scientific support, and the remaining omissions limit broader claims.
A live comparison needs a fixed model version, held-out questions, matched
information, independent adjudication and a budget. \citep{rel66pharma,rel66inferbert,rel66deepcausality,rel66veitch}





